\documentclass[12pt,a4paper]{article}

\usepackage[a4paper,margin=2cm]{geometry}
\usepackage[T2A]{fontenc}
\usepackage[utf8]{inputenc}
\usepackage[russian]{babel}
\usepackage{amsmath,amssymb,mathtools}
\usepackage{array,booktabs,longtable}
\usepackage{xcolor}
\usepackage{hyperref}
\usepackage{tikz}

\hypersetup{
  colorlinks=true,
  linkcolor=blue!55!black,
  urlcolor=blue!55!black
}

\setlength{\parindent}{0pt}
\setlength{\parskip}{0.55em}
\renewcommand{\arraystretch}{1.22}

\begin{document}

\begin{titlepage}
\centering
\vspace*{0.24\textheight}
{\LARGE\bfseries Ревью домашней работы\par}
\vspace{1.6cm}
{\large Автор: Лёвин Артём Александрович\par}
\vspace{0.8cm}
{\large 16 сентября 2026 года\par}
\vfill
\begin{tikzpicture}
  \draw[blue!35!black, line width=0.8pt]
    (0,0) .. controls (2,0.45) and (4,-0.45) .. (6,0)
          .. controls (8,0.45) and (10,-0.45) .. (12,0);
\end{tikzpicture}
\end{titlepage}

\newpage
\section*{Сводная таблица}
\small
\begin{center}
\begin{tabular}{>{\centering\arraybackslash}p{0.08\linewidth}
                p{0.18\linewidth}
                p{0.29\linewidth}
                p{0.14\linewidth}
                p{0.14\linewidth}}
\toprule
№ задания & Тема & Суть ошибки & Ваш ответ & Правильный ответ \\
\midrule
\hyperref[task:2]{2} & Попытки до первого успеха & После успеха дерево продолжено, поэтому в сумму попали ветви, которых по условию уже не существует. & $0{,}8$ & $0{,}98$ \\
\addlinespace
\hyperref[task:8]{8} & Две игры и сумма очков & Сложены исходы, дающие меньше четырёх очков, то есть найдено дополнение к нужному событию. & $0{,}55$ & $0{,}45$ \\
\addlinespace
\hyperref[task:10]{10} & Полная вероятность & Дерево начато верно, но вычисление не доведено до итоговой вероятности; окончательный ответ не записан. & не указан & $0{,}846$ \\
\bottomrule
\end{tabular}
\end{center}
\normalsize

\newpage
\vspace*{\fill}
\begin{center}
{\LARGE\bfseries Разбор ошибок}
\end{center}
\vspace*{\fill}

\newpage
\phantomsection\label{task:2}
\section*{Задание 2}

\textbf{Текст задания.}
Вероятность успеха при первой попытке равна $0{,}5$, а при каждой следующей --- $0{,}8$. Попытки продолжаются до первого успеха. Найдите вероятность того, что успех произойдёт не позднее третьей попытки.

\textbf{Ваше решение.}
Построено дерево исходов для трёх попыток. После первой попытки указаны вероятности успеха и неуспеха $0{,}5$ и $0{,}5$, а для следующих попыток --- $0{,}8$ и $0{,}2$. Однако ветвление продолжается и после тех ветвей, на которых успех уже произошёл. Затем складываются вероятности нескольких трёхшаговых ветвей.

\textbf{Ваш ответ.} $0{,}8$.

\textcolor{red}{\textbf{Ошибка:}} опыт продолжен после первого успеха.

\textbf{Где именно возникает ошибка.}
До первого успеха дерево строится правильно: на первой попытке есть успех с вероятностью $0{,}5$ и неуспех с вероятностью $0{,}5$; после первого неуспеха на второй попытке есть успех с вероятностью $0{,}8$ и неуспех с вероятностью $0{,}2$.

Первый неверный шаг появляется тогда, когда после ветви <<успех>> снова строится следующая попытка. По условию после первого успеха попытки прекращаются, поэтому таких продолжений у соответствующего исхода быть не должно.

\textbf{Почему этот шаг неверен.}
Событие <<успех не позднее третьей попытки>> состоит из трёх несовместимых вариантов:
\begin{enumerate}
  \item успех сразу на первой попытке;
  \item первая попытка неуспешна, а вторая успешна;
  \item первые две попытки неуспешны, а третья успешна.
\end{enumerate}
После того как успех произошёл, дальнейших попыток уже нет. Поэтому нельзя добавлять к ветви успеха новые множители и затем считать такие продолжения отдельными исходами: это меняет сам эксперимент, заданный условием.

\textcolor{blue!60!black}{\textbf{Ключевая идея:}} дерево должно обрываться на первом успехе. Тогда достаточно сложить вероятности трёх способов впервые получить успех не позднее третьей попытки.

\textbf{Исправление от последнего правильного шага.}
Оставим только допустимые ветви:
\[
P_1=0{,}5,
\]
\[
P_2=0{,}5\cdot 0{,}8=0{,}4,
\]
\[
P_3=0{,}5\cdot 0{,}2\cdot 0{,}8=0{,}08.
\]
Складываем, потому что эти три случая не могут произойти одновременно:
\[
P=P_1+P_2+P_3=0{,}5+0{,}4+0{,}08=0{,}98.
\]

\textbf{Полное правильное решение.}
Проще всего рассматривать момент первого успеха.

На первой попытке успех происходит с вероятностью $0{,}5$.

Чтобы первый успех произошёл на второй попытке, первая должна закончиться неуспехом, а вторая --- успехом. Вероятность этого равна
\[
0{,}5\cdot 0{,}8=0{,}4.
\]

Чтобы первый успех произошёл на третьей попытке, первые две попытки должны закончиться неуспехом, а третья --- успехом. Вероятность равна
\[
0{,}5\cdot 0{,}2\cdot 0{,}8=0{,}08.
\]

Следовательно,
\[
P=0{,}5+0{,}4+0{,}08=0{,}98.
\]

\textbf{Проверка.}
Можно проверить ответ через противоположное событие: успех не произошёл ни в одной из первых трёх попыток. Для этого первая попытка должна быть неуспешной с вероятностью $0{,}5$, а вторая и третья --- с вероятностью $0{,}2$ каждая:
\[
P(\text{три неуспеха})=0{,}5\cdot0{,}2\cdot0{,}2=0{,}02.
\]
Тогда
\[
1-0{,}02=0{,}98,
\]
что совпадает с полученным результатом.

\textcolor{teal!60!black}{\textbf{Итог:}} вероятность того, что успех произойдёт не позднее третьей попытки, равна $0{,}98$.

\textbf{Задача на закрепление.}
Вероятность успеха при первой попытке равна $0{,}6$, а при каждой следующей --- $0{,}75$. Попытки продолжаются до первого успеха. Найдите вероятность того, что успех произойдёт не позднее третьей попытки.

\newpage
\phantomsection\label{task:8}
\section*{Задание 8}

\textbf{Текст задания.}
Футбольная команда играет две независимые игры. В каждой игре вероятность победы равна $0{,}5$, ничьей --- $0{,}2$, поражения --- $0{,}3$. За победу начисляют $3$ очка, за ничью --- $1$, за поражение --- $0$. Найдите вероятность того, что за две игры команда наберёт не менее $4$ очков.

\textbf{Ваше решение.}
Построено дерево исходов с вероятностями $0{,}5$, $0{,}2$ и $0{,}3$ для победы, ничьей и поражения в каждой игре. В итоговой сумме записаны вероятности шести комбинаций и получено
\[
0{,}55.
\]

\textbf{Ваш ответ.} $0{,}55$.

\textcolor{red}{\textbf{Ошибка:}} выбраны ветви противоположного события.

\textbf{Где именно возникает ошибка.}
Само дерево вероятностей построено по условию корректно. Ошибка начинается при выборе комбинаций, которые нужно складывать. В сумму попали исходы, дающие меньше $4$ очков: например, победа и поражение дают $3$ очка, две ничьи дают $2$ очка, ничья и поражение дают $1$ очко.

Именно эти шесть исходов в совокупности имеют вероятность $0{,}55$. Это вероятность события <<набрать меньше $4$ очков>>, а не события из условия.

\textbf{Почему этот шаг неверен.}
Фраза <<не менее $4$ очков>> означает $4$ очка или больше. За две игры это возможно только в трёх случаях:
\begin{itemize}
  \item победа и победа: $6$ очков;
  \item победа и ничья: $4$ очка;
  \item ничья и победа: $4$ очка.
\end{itemize}
Все остальные комбинации дают $3$, $2$, $1$ или $0$ очков и потому не удовлетворяют условию.

\textcolor{blue!60!black}{\textbf{Ключевая идея:}} сначала определить, какие пары результатов дают нужное число очков, и только после этого складывать вероятности соответствующих ветвей.

\textbf{Исправление от последнего правильного шага.}
Из уже построенного дерева оставляем только три подходящие ветви:
\[
P(\text{победа, победа})=0{,}5\cdot0{,}5=0{,}25,
\]
\[
P(\text{победа, ничья})=0{,}5\cdot0{,}2=0{,}10,
\]
\[
P(\text{ничья, победа})=0{,}2\cdot0{,}5=0{,}10.
\]
Складываем:
\[
P=0{,}25+0{,}10+0{,}10=0{,}45.
\]

\textbf{Полное правильное решение.}
Перечислим только те исходы двух игр, при которых команда получает не менее $4$ очков.

Две победы дают $3+3=6$ очков. Вероятность:
\[
0{,}5\cdot0{,}5=0{,}25.
\]

Победа в первой игре и ничья во второй дают $3+1=4$ очка. Вероятность:
\[
0{,}5\cdot0{,}2=0{,}10.
\]

Ничья в первой игре и победа во второй также дают $1+3=4$ очка. Вероятность:
\[
0{,}2\cdot0{,}5=0{,}10.
\]

Эти исходы несовместимы, поэтому
\[
P=0{,}25+0{,}10+0{,}10=0{,}45.
\]

\textbf{Проверка.}
В Вашем вычислении получено $0{,}55$ для остальных шести исходов, то есть для события <<набрать меньше $4$ очков>>. Поэтому
\[
0{,}45+0{,}55=1.
\]
Сумма равна единице, что подтверждает, что два события являются противоположными и правильная вероятность равна $0{,}45$.

\textcolor{teal!60!black}{\textbf{Итог:}} вероятность набрать за две игры не менее $4$ очков равна $0{,}45$.

\textbf{Задача на закрепление.}
Команда играет две независимые игры. В каждой игре вероятность победы равна $0{,}4$, ничьей --- $0{,}3$, поражения --- $0{,}3$. За победу начисляют $3$ очка, за ничью --- $1$, за поражение --- $0$. Найдите вероятность того, что за две игры команда наберёт не менее $4$ очков.

\newpage
\phantomsection\label{task:10}
\section*{Задание 10}

\textbf{Текст задания.}
Стрелок случайно выбирает одно из двух устройств и после выбора делает из него два независимых выстрела. Первое устройство выбирается с вероятностью $0{,}4$, второе --- с вероятностью $0{,}6$. Вероятность попадания одним выстрелом из первого устройства равна $0{,}9$, а из второго --- $0{,}5$. Найдите вероятность того, что хотя бы один из двух выстрелов попадёт в цель.

\textbf{Ваше решение.}
Начато дерево вероятностей: выбор первого устройства с вероятностью $0{,}4$ и второго с вероятностью $0{,}6$ указан верно; для первого устройства используются вероятности попадания и промаха $0{,}9$ и $0{,}1$, для второго --- $0{,}5$ и $0{,}5$. Далее вычисление перечёркнуто и завершённого значения вероятности нет.

\textbf{Ваш ответ.} не указан.

\textcolor{red}{\textbf{Ошибка:}} решение не доведено до итогового ответа.

\textbf{Где именно возникает ошибка.}
Последний надёжно выполненный шаг --- разбиение на два случая по выбранному устройству и указание вероятностей попадания и промаха для каждого устройства.

Первый неполный шаг --- после построения ветвей не вычислена суммарная вероятность нужного события. Из-за этого решение обрывается до получения численного ответа.

\textbf{Почему это важно.}
Вероятность попадания зависит от того, какое устройство выбрано. Поэтому сначала нужно учесть вероятность выбора устройства, а затем --- вероятность двух выстрелов при этом выборе. Простое сложение вероятностей попадания отдельных выстрелов здесь не подходит: события пересекаются, а устройство выбирается один раз перед двумя выстрелами.

Удобнее считать противоположное событие: оба выстрела промахнулись. Оно легко вычисляется отдельно для каждого устройства, после чего применяется формула полной вероятности.

\textcolor{blue!60!black}{\textbf{Ключевая идея:}} найти вероятность двух промахов с учётом выбора устройства, а затем вычесть её из единицы.

\textbf{Исправление от последнего правильного шага.}
Если выбрано первое устройство, вероятность промаха одним выстрелом равна
\[
1-0{,}9=0{,}1.
\]
Два выстрела независимы, поэтому вероятность двух промахов равна
\[
0{,}1\cdot0{,}1=0{,}01.
\]
С учётом вероятности выбора первого устройства получаем вклад
\[
0{,}4\cdot0{,}01=0{,}004.
\]

Если выбрано второе устройство, вероятность промаха одним выстрелом равна
\[
1-0{,}5=0{,}5,
\]
а вероятность двух промахов
\[
0{,}5\cdot0{,}5=0{,}25.
\]
С учётом выбора второго устройства вклад равен
\[
0{,}6\cdot0{,}25=0{,}15.
\]

Значит, вероятность двух промахов после случайного выбора устройства:
\[
0{,}004+0{,}15=0{,}154.
\]
Тогда вероятность хотя бы одного попадания:
\[
1-0{,}154=0{,}846.
\]

\textbf{Полное правильное решение.}
Рассмотрим противоположное событие: оба выстрела оказались промахами.

Для первого устройства вероятность двух промахов равна $0{,}1^2=0{,}01$, а само устройство выбирается с вероятностью $0{,}4$. Поэтому вероятность получить первый случай и два промаха равна
\[
0{,}4\cdot0{,}01=0{,}004.
\]

Для второго устройства вероятность двух промахов равна $0{,}5^2=0{,}25$, а устройство выбирается с вероятностью $0{,}6$. Поэтому соответствующая вероятность равна
\[
0{,}6\cdot0{,}25=0{,}15.
\]

Эти два случая несовместимы, следовательно,
\[
P(\text{два промаха})=0{,}004+0{,}15=0{,}154.
\]
Искомая вероятность:
\[
P(\text{хотя бы одно попадание})=1-0{,}154=0{,}846.
\]

\textbf{Проверка.}
Проверим тем же результатом через вероятность хотя бы одного попадания при каждом устройстве.

Для первого устройства:
\[
1-0{,}1^2=0{,}99.
\]
Для второго устройства:
\[
1-0{,}5^2=0{,}75.
\]
По формуле полной вероятности:
\[
0{,}4\cdot0{,}99+0{,}6\cdot0{,}75=0{,}396+0{,}45=0{,}846.
\]
Результат совпал.

\textcolor{teal!60!black}{\textbf{Итог:}} вероятность того, что хотя бы один из двух выстрелов попадёт в цель, равна $0{,}846$.

\textbf{Задача на закрепление.}
Стрелок выбирает одно из двух устройств: первое с вероятностью $0{,}3$, второе с вероятностью $0{,}7$. После выбора он делает два независимых выстрела. Вероятность попадания одним выстрелом из первого устройства равна $0{,}8$, а из второго --- $0{,}6$. Найдите вероятность того, что хотя бы один из двух выстрелов попадёт в цель.

\vfill
\begin{center}
\small Лёвин Артём Александрович эксклюзивно для Ксении Клыковой
\end{center}

\end{document}
